\documentclass[11pt,a4paper]{article}

% ── Packages ──
\usepackage[utf8]{inputenc}
\usepackage[T1]{fontenc}
\usepackage{lmodern}
\usepackage[margin=2.4cm]{geometry}
\usepackage{amsmath,amssymb,amsfonts,mathtools}
\usepackage{bm}
\usepackage{booktabs}
\usepackage{enumitem}
\usepackage{hyperref}
\usepackage{cleveref}
\usepackage{xcolor}
\usepackage{tcolorbox}
\usepackage{tikz}
\usetikzlibrary{arrows.meta,positioning,calc,fit,backgrounds}
\usepackage{fancyhdr}
\usepackage{titlesec}
\usepackage{parskip}

% ── Colors ──
\definecolor{boxblue}{RGB}{230,240,255}
\definecolor{boxborder}{RGB}{70,130,200}
\definecolor{examplegreen}{RGB}{235,250,235}
\definecolor{exampleborder}{RGB}{60,160,60}
\definecolor{reasoningamber}{RGB}{255,248,230}
\definecolor{reasoningborder}{RGB}{200,160,40}

% ── Custom boxes ──
\newtcolorbox{keqbox}[1][]{
  colback=boxblue, colframe=boxborder, arc=3pt,
  boxrule=0.8pt, left=6pt, right=6pt, top=4pt, bottom=4pt,
  title={#1}, fonttitle=\bfseries\sffamily
}

\newtcolorbox{examplebox}[1][Example]{
  colback=examplegreen, colframe=exampleborder, arc=3pt,
  boxrule=0.6pt, left=6pt, right=6pt, top=4pt, bottom=4pt,
  title={#1}, fonttitle=\bfseries\sffamily\small
}

\newtcolorbox{reasonbox}[1][Reasoning]{
  colback=reasoningamber, colframe=reasoningborder, arc=3pt,
  boxrule=0.6pt, left=6pt, right=6pt, top=4pt, bottom=4pt,
  title={#1}, fonttitle=\bfseries\sffamily\small
}

% ── Shortcuts ──
\newcommand{\R}{\mathbb{R}}
\newcommand{\SE}{\mathrm{SE}}
\newcommand{\SO}{\mathrm{SO}}
\newcommand{\so}{\mathfrak{so}}
\newcommand{\se}{\mathfrak{se}}
\newcommand{\Exp}{\operatorname{Exp}}
\newcommand{\Log}{\operatorname{Log}}
\newcommand{\tr}{\mathsf{T}}
\newcommand{\diag}{\operatorname{diag}}
\newcommand{\clamp}{\operatorname{clamp}}
\newcommand{\SLERP}{\operatorname{SLERP}}
\newcommand{\odot}{\circ}  % element-wise product

% ── Headers ──
\pagestyle{fancy}
\fancyhf{}
\fancyhead[L]{\small\sffamily AIC Controller --- Mathematical Reference}
\fancyhead[R]{\small\sffamily\thepage}
\renewcommand{\headrulewidth}{0.4pt}

% ── Section formatting ──
\titleformat{\section}{\Large\bfseries\sffamily}{{\thesection}}{0.8em}{}
\titleformat{\subsection}{\large\bfseries\sffamily}{{\thesubsection}}{0.6em}{}
\titleformat{\subsubsection}{\normalsize\bfseries\sffamily}{{\thesubsubsection}}{0.5em}{}

\hypersetup{
  colorlinks=true, linkcolor=boxborder, urlcolor=boxborder, citecolor=boxborder
}

% ══════════════════════════════════════════════════════════════════════════════
\begin{document}

\begin{center}
  {\LARGE\bfseries\sffamily AIC Controller}\\[4pt]
  {\Large\sffamily Mathematical Foundations \& Algorithm Reference}\\[12pt]
  {\normalsize Auto-generated from source --- \texttt{aic\_controller} package}\\[2pt]
  {\small\today}
\end{center}

\vspace{6pt}
\hrule height 0.8pt
\vspace{10pt}

\tableofcontents
\newpage

% ══════════════════════════════════════════════════════════════════════════════
\section{Architecture Overview}
\label{sec:arch}

The AIC Controller is a ROS\,2 \texttt{controller\_interface::ControllerInterface}
plugin.  It runs inside \texttt{ros2\_control} at a fixed frequency and supports
two \emph{control modes}:

\begin{center}
\begin{tabular}{lll}
\toprule
\textbf{Mode} & \textbf{Output Interface} & \textbf{Status} \\
\midrule
Impedance  & Joint effort (torque) & Implemented \\
Admittance & Joint position        & Not yet implemented \\
\bottomrule
\end{tabular}
\end{center}

Within impedance mode, two \emph{target modes} select the coordinate space of
incoming commands:

\begin{center}
\begin{tabular}{lll}
\toprule
\textbf{Target Mode} & \textbf{Message Type} & \textbf{Impedance Action} \\
\midrule
Cartesian & \texttt{MotionUpdate}      & \texttt{CartesianImpedanceAction} \\
Joint     & \texttt{JointMotionUpdate} & \texttt{JointImpedanceAction} \\
\bottomrule
\end{tabular}
\end{center}

Each control cycle executes:

\begin{center}
\fbox{\textsf{Read Hardware}}\;$\to$\;
\fbox{\textsf{Ingest Commands}}\;$\to$\;
\fbox{\textsf{Generate Reference}}\;$\to$\;
\fbox{\textsf{Compute Torques}}\;$\to$\;
\fbox{\textsf{Write Hardware}}
\end{center}


% ══════════════════════════════════════════════════════════════════════════════
\section{Mathematical Preliminaries}
\label{sec:prelim}

% --------------------------------------------------------------------------
\subsection{Rigid Body Poses --- $\SE(3)$ and $\SO(3)$}
\label{sec:SE3}

A rigid body pose in 3D is an element of the \emph{Special Euclidean group}:
\[
  T = \begin{bmatrix} R & p \\ \bm{0}^\tr & 1 \end{bmatrix} \in \SE(3),
  \qquad R \in \SO(3),\; p \in \R^3
\]
where $R$ is a rotation matrix and $p$ is a translation vector.

In code, poses are stored as \texttt{Eigen::Isometry3d} ---
\texttt{pose.linear()} gives $R$ and \texttt{pose.translation()} gives $p$.

% --------------------------------------------------------------------------
\subsection{Exponential and Logarithmic Maps}
\label{sec:explog}

The Lie algebra $\so(3)$ consists of $3\!\times\!3$ skew-symmetric matrices.
The \emph{hat} operator maps $\R^3 \to \so(3)$:
\[
  \bm{\omega}
  = \begin{bmatrix} \omega_1 \\ \omega_2 \\ \omega_3 \end{bmatrix}
  \;\mapsto\;
  [\bm{\omega}]_\times
  = \begin{bmatrix}
      0         & -\omega_3 & \omega_2 \\
      \omega_3  & 0         & -\omega_1 \\
      -\omega_2 & \omega_1  & 0
    \end{bmatrix}
\]

\paragraph{Exponential map} (tangent vector $\to$ rotation):
\[
  \Exp : \R^3 \to \SO(3),
  \qquad
  R = \exp\!\bigl([\bm{\omega}]_\times\bigr)
\]

Computed via the Rodrigues formula.  For quaternions (Sophus \texttt{SO3::exp}):
\[
  q = \Exp(\bm{\delta}),
  \qquad
  \theta = \|\bm{\delta}\|,
  \quad
  q = \Bigl(\cos\tfrac{\theta}{2},\;
             \tfrac{\bm{\delta}}{\theta}\sin\tfrac{\theta}{2}\Bigr)
\]

\paragraph{Logarithmic map} (rotation $\to$ tangent vector):
\[
  \Log : \SO(3) \to \R^3,
  \qquad
  \bm{\omega} = \Log(R)
\]
This is the inverse of $\Exp$.  Given a unit quaternion $q$, it extracts the
axis-angle vector $\bm{\omega}$ with $\|\bm{\omega}\| = \theta$ (rotation angle)
and $\bm{\omega}/\theta$ the rotation axis.

\paragraph{SE(3) integration.}
The twist is $\bm{\xi} = [v^\tr,\, \omega^\tr]^\tr \in \R^6$.
Pose integration uses:
\begin{keqbox}[SE(3) Integration]
\[
  T_{k+1} = T_k \cdot \Exp(\bm{\xi}\,\Delta t)
\]
\end{keqbox}

Implemented in \texttt{utils::integrate\_pose} via Sophus \texttt{SE3::exp}.

\begin{reasonbox}[Why Lie groups instead of Euler angles?]
Classical Euler angles suffer from gimbal lock and representation singularities.
Working directly on $\SE(3)$/$\SO(3)$ ensures singularity-free pose integration
and consistent orientation error computation.
\end{reasonbox}

% --------------------------------------------------------------------------
\subsection{Robot Jacobian}
\label{sec:jacobian}

For an $n$-DOF robot, the geometric Jacobian $J(q) \in \R^{6\times n}$
maps joint velocities to the end-effector twist:
\[
  \bm{\xi} = J(q)\,\dot{q}
\]
where $\bm{\xi} = [v^\tr,\,\omega^\tr]^\tr$ is the 6D spatial velocity
(linear + angular) of the tool frame.

The Jacobian is computed externally by a \texttt{kinematics\_interface} plugin
(typically KDL- or Pinocchio-based) via \texttt{calculate\_jacobian()}.


% ══════════════════════════════════════════════════════════════════════════════
\section{Cartesian Impedance Control}
\label{sec:cartesian_imp}

% --------------------------------------------------------------------------
\subsection{The Spring--Damper Analogy}

Impedance control makes the robot behave like a mechanical spring--damper system
in Cartesian space.  A~1D mass--spring--damper satisfies:
\[
  m\ddot{x} + d\dot{x} + kx = f_{\text{ext}}
\]

In impedance control we \emph{generate} the wrench:
\[
  F = K \cdot e_x + D \cdot e_v
\]
where $e_x$ is position error and $e_v$ is velocity error, allowing compliant
interaction with the environment.

% --------------------------------------------------------------------------
\subsection{Control Wrench Law}
\label{sec:wrench_law}

The full 6D control wrench in the base frame is:

\begin{keqbox}[Cartesian Impedance Wrench]
\begin{equation}
  F_{\text{ctrl}}
  = K\,e_x + D\,e_v + F_{\text{ff}} + K_I \odot \!\int\! e_x\,dt + F_{\text{offset}}
  \label{eq:wrench}
\end{equation}
\end{keqbox}

\begin{center}
\begin{tabular}{ccc}
\toprule
\textbf{Symbol} & \textbf{Size} & \textbf{Description} \\
\midrule
$K$ & $6\!\times\!6$ & Cartesian stiffness (diagonal) \\
$D$ & $6\!\times\!6$ & Cartesian damping (diagonal) \\
$e_x$ & $6\!\times\!1$ & Pose error $= x_{\text{des}} - x_{\text{cur}}$ \\
$e_v$ & $6\!\times\!1$ & Velocity error $= \dot{x}_{\text{des}} - \dot{x}_{\text{cur}}$ \\
$F_{\text{ff}}$ & $6\!\times\!1$ & Feedforward wrench (see \S\ref{sec:smoothing}) \\
$K_I$ & $6\!\times\!1$ & Integrator gain (element-wise) \\
$F_{\text{offset}}$ & $6\!\times\!1$ & Constant offset (e.g.\ payload) \\
\bottomrule
\end{tabular}
\end{center}

The pose error $e_x$ is computed by the kinematics plugin's
\texttt{calculate\_frame\_difference()}, yielding
$[e_{\text{trans}}^\tr,\, e_{\text{rot}}^\tr]^\tr$ where the rotational part is
in axis-angle form.

\begin{examplebox}
Suppose TCP position $p = [0.5,\,0,\,0.3]$\,m and desired
$p_d = [0.5,\,0.05,\,0.3]$\,m with
$K = \diag(500,500,500,20,20,20)$\,N/m\,(N\,m/rad):
\[
  e_x = [0,\; 0.05,\; 0,\; 0,\; 0,\; 0]^\tr
  \;\;\Longrightarrow\;\;
  F = K\,e_x = [0,\; 25,\; 0,\; 0,\; 0,\; 0]^\tr\;\text{N}
\]
A 25\,N restoring force along $y$.
\end{examplebox}

% --------------------------------------------------------------------------
\subsection{PID-Like Extension --- Integral Term}
\label{sec:integral}

An integral term accumulates pose error to eliminate steady-state offset:
\begin{align}
  e_I(k\!+\!1) &= \clamp\!\bigl(e_I(k) + e_x(k),\; -B,\; B\bigr)
  \label{eq:integrator_update} \\[4pt]
  F_I &= K_I \odot e_I
  \label{eq:integrator_force}
\end{align}

where $B \in \R^6$ is the integrator bound (anti-windup) and $\odot$ denotes
element-wise multiplication.

\begin{reasonbox}[Why anti-windup?]
Without the bound $B$, large transient errors (e.g.\ during collisions) cause
the integrator to \emph{wind up}, producing dangerous persistent forces even
after the error is resolved.
\end{reasonbox}

% --------------------------------------------------------------------------
\subsection{Mapping Wrench to Joint Torques}
\label{sec:JT}

The Cartesian wrench is mapped to joint torques via the Jacobian transpose:

\begin{keqbox}[Jacobian Transpose Mapping]
\begin{equation}
  \tau = J^\tr\, F_{\text{ctrl}}
  \label{eq:JT}
\end{equation}
\end{keqbox}

\paragraph{Derivation from virtual work.}
A virtual joint displacement $\delta q$ produces a Cartesian displacement
$\delta x = J\,\delta q$.  The work done by force $F$ is:
\[
  \delta W
  = F^\tr \delta x
  = F^\tr J\,\delta q
  = (J^\tr F)^\tr \delta q
\]
So the equivalent joint torque is $\tau = J^\tr F$.  This does \emph{not}
require inverting the Jacobian and works even at singularities (though
Cartesian stiffness degenerates in singular directions).

% --------------------------------------------------------------------------
\subsection{Wrench and Torque Saturation}

Before applying $J^\tr$, each wrench component is clamped:
\[
  F_{\text{ctrl},i}
  = \clamp(F_{\text{ctrl},i},\; -F_{\max,i},\; F_{\max,i})
\]

After computing all joint torque contributions, each joint torque is clamped to
the URDF effort limit:
\[
  \tau_i = \clamp(\tau_i,\; -\tau_{\max,i},\; \tau_{\max,i})
\]


% ══════════════════════════════════════════════════════════════════════════════
\section{Nullspace Control}
\label{sec:nullspace}

% --------------------------------------------------------------------------
\subsection{Redundancy and the Nullspace}

A robot with $n > 6$ joints has \emph{kinematic redundancy}.  The extra DOFs
span the nullspace of $J$:
\[
  \mathcal{N}(J)
  = \bigl\{\, \dot{q} \in \R^n \mid J\,\dot{q} = \bm{0} \,\bigr\}
\]

Nullspace control adds torques that move joints \emph{without} affecting the TCP,
useful for preferred-posture regulation, joint-limit avoidance, or manipulability
optimisation.

% --------------------------------------------------------------------------
\subsection{Nullspace Projector}

The nullspace projection matrix is:
\begin{equation}
  N = I_n - J\,J^\dagger
  \label{eq:nullspace_proj}
\end{equation}
where $J^\dagger$ is the right pseudo-inverse of $J$.  Any torque
$\tau_0 \in \R^n$ projected through $N$ produces zero TCP acceleration:
\[
  \tau_{\text{ns}} = N\,\tau_0
\]

\paragraph{Proof sketch.}
$J\,J^\dagger\,J = J$, so $J\,N = J(I - JJ^\dagger) = J - J = \bm{0}$.
Therefore the Cartesian acceleration produced by $\tau_{\text{ns}}$ is zero in
the linearised dynamics.

% --------------------------------------------------------------------------
\subsection{Singularity-Robust Pseudo-Inverse (SVD Regularisation)}
\label{sec:svd}

The na\"ive pseudo-inverse $J^\dagger = J^\tr(JJ^\tr)^{-1}$ diverges near
singularities.  This implementation uses a \emph{damped pseudo-inverse}
following \textbf{Chiaverini (1997)} with per-dimension adaptive damping.

\medskip

Compute the SVD of $M = JJ^\tr$:
\[
  M = U\,\Sigma\,V^\tr,
  \qquad
  \Sigma = \diag(\sigma_1,\, \sigma_2,\, \ldots,\, \sigma_m),
  \quad
  \sigma_1 \ge \sigma_2 \ge \cdots \ge \sigma_m \ge 0
\]

The regularised inverse diagonal entries are:

\begin{keqbox}[Regularised Singular Values]
\begin{equation}
  \Sigma^\dagger_{ii}
  = \frac{\sigma_i}{\sigma_i^2 + \lambda_i^2}
  \label{eq:reg_sv}
\end{equation}
\end{keqbox}

The per-dimension condition number is:
\[
  \text{cn}_i = \frac{\sigma_1}{\sigma_i + 10^{-10}}
\]

The adaptive damping factor is:
\begin{equation}
  \lambda_i^2 =
  \begin{cases}
    \displaystyle
    \left(1 - \Bigl(\frac{c_{\text{th}}}{\text{cn}_i}\Bigr)^{\!2}\right)^{\!2}
    \cdot \frac{1}{\epsilon}
    & \text{if } \text{cn}_i > c_{\text{th}} \\[10pt]
    0 & \text{otherwise}
  \end{cases}
  \label{eq:lambda}
\end{equation}

where $c_{\text{th}} = 1500$ is the condition number threshold and
$\epsilon = 10^{-6}$.

\paragraph{The bisquare function $f(r) = (1 - r^2)^2$.}
\begin{itemize}[nosep]
  \item $f(0) = 1$: full damping when condition number is extremely high
  \item $f(1) = 0$: zero damping at the threshold boundary
  \item $f'(0) = 0$ and $f'(1) = 0$: smooth onset/offset --- no discontinuous jumps
\end{itemize}

The final pseudo-inverse is:
\begin{equation}
  J^\dagger = J^\tr \, V \, \Sigma^\dagger \, U^\tr
  \label{eq:pinv}
\end{equation}

\begin{examplebox}[Singularity example]
Consider a 7-DOF arm near a shoulder singularity where
$\sigma_6 \approx 0.001$ and $\sigma_1 = 5.0$:
\[
  \text{cn}_6 = \frac{5.0}{0.001} = 5000 \;\gg\; 1500
  \qquad\Longrightarrow\qquad
  \lambda_6^2 = \left(1 - (1500/5000)^2\right)^2 \cdot 10^6
  \approx 7.6\!\times\!10^5
\]
This large $\lambda_6^2$ suppresses the near-zero singular value, preventing
torque spikes.
\end{examplebox}

% --------------------------------------------------------------------------
\subsection{Nullspace Impedance Torque}

The raw nullspace torque is a joint-space PD controller:
\begin{equation}
  \tau_0
  = K_{\text{ns}} \odot (q_{\text{ns}}^* - q) - D_{\text{ns}} \odot \dot{q}
  \label{eq:ns_raw}
\end{equation}

where $q_{\text{ns}}^*$ is the desired nullspace configuration and
$K_{\text{ns}},\,D_{\text{ns}} \in \R^n$ are per-joint stiffness/damping.

Projected through the nullspace:
\begin{equation}
  \tau_{\text{ns}} = N\,\tau_0
  = (I - J\,J^\dagger)\,\tau_0
  \label{eq:ns_projected}
\end{equation}


% ══════════════════════════════════════════════════════════════════════════════
\section{Joint-Limit Avoidance}
\label{sec:jlimit}

% --------------------------------------------------------------------------
\subsection{Activation Zone}

For each joint $k$, an activation zone is defined as a percentage $\alpha\in(0,1]$
of the joint range:
\begin{align}
  q_{\text{range}} &= q_{\max} - q_{\min} \notag\\[4pt]
  q_{\text{upper\_th}} &= q_{\min} + \frac{\alpha\, q_{\text{range}}}{2}
  \label{eq:upper_thresh} \\[4pt]
  q_{\text{lower\_th}} &= q_{\max} - \frac{\alpha\, q_{\text{range}}}{2}
  \label{eq:lower_thresh}
\end{align}

\begin{center}
\begin{tikzpicture}[>=Stealth, thick, font=\small\sffamily]
  \draw[->] (0,0) -- (10.5,0) node[right]{$q$};
  \foreach \x/\lab in {0.5/$q_{\min}$, 3/$q_{\text{lower\_th}}$,
                        7/$q_{\text{upper\_th}}$, 9.5/$q_{\max}$}{
    \draw (\x,0.15) -- (\x,-0.15) node[below]{\lab};
  }
  \draw[red!70!black, very thick] (0.5,0) -- (3,0);
  \draw[green!50!black, very thick] (3,0) -- (7,0);
  \draw[red!70!black, very thick] (7,0) -- (9.5,0);
  \node[red!70!black, above] at (1.75,0.2) {active};
  \node[green!50!black, above] at (5,0.2) {safe zone};
  \node[red!70!black, above] at (8.25,0.2) {active};
\end{tikzpicture}
\end{center}

% --------------------------------------------------------------------------
\subsection{Potential Field Torque}

Within the activation zone, a linear repulsive potential pushes the joint away
from its limit:

\begin{equation}
  \tau_{\text{avoid},k} =
  \begin{cases}
    g_u \cdot (q_{\text{upper\_th}} - q_k) & \text{if } q_k > q_{\text{upper\_th}} \\[6pt]
    g_l \cdot (q_{\text{lower\_th}} - q_k) & \text{if } q_k < q_{\text{lower\_th}} \\[6pt]
    0 & \text{otherwise}
  \end{cases}
  \label{eq:avoid}
\end{equation}

The gains are chosen so the torque reaches $\tau_{\max}$ exactly at the limit:
\begin{equation}
  g_u = \frac{\tau_{\max}}{|q_{\max} - q_{\text{upper\_th}}|},
  \qquad
  g_l = \frac{\tau_{\max}}{|q_{\min} - q_{\text{lower\_th}}|}
  \label{eq:avoid_gains}
\end{equation}

If the joint exceeds its limit, the torque is \emph{saturated} at $\tau_{\max}$.


% ══════════════════════════════════════════════════════════════════════════════
\section{Joint-Space Impedance Control}
\label{sec:joint_imp}

For joint-target mode, a diagonal impedance law is used:

\begin{keqbox}[Joint Impedance Law]
\begin{equation}
  \tau
  = K \odot (q^* - q) + D \odot (\dot{q}^* - \dot{q}) + \tau_{\text{ff}}
  \label{eq:joint_imp}
\end{equation}
\end{keqbox}

where $K,\,D \in \R^n$ are per-joint stiffness/damping vectors,
$q^*,\,\dot{q}^*$ are reference position/velocity, $\tau_{\text{ff}}$ is
feedforward torque, and all operations are element-wise.

Each torque is clamped: $\tau_i = \clamp(\tau_i,\,-\tau_{\max,i},\,\tau_{\max,i})$.

This mode is simpler than Cartesian impedance --- it operates directly in joint
space, requiring no Jacobian, no nullspace projection, and no FK-based error
computation.


% ══════════════════════════════════════════════════════════════════════════════
\section{Gravity Compensation}
\label{sec:gravity}

The gravity torque vector $\tau_g(q)$ compensates for gravitational forces on
robot links:
\begin{equation}
  \tau_g(q)
  = \sum_{i=1}^{n} J_{v,i}^\tr(q)\, m_i\, \mathbf{g}
  \label{eq:gravity}
\end{equation}

where $J_{v,i}$ is the linear Jacobian of the CoM of link $i$, $m_i$ is its
mass, and $\mathbf{g} = [0,\, 0,\, -9.80665]^\tr$\,m/s$^2$.

Computed by KDL's \texttt{ChainDynParam::JntToGravity} from the URDF kinematic
chain.  The result is \emph{added} to the impedance torques:
\[
  \tau_{\text{total}} = \tau_{\text{impedance}} + \tau_g(q)
\]

\begin{reasonbox}
On real hardware, motor controllers receive raw torque commands, so the
controller must explicitly cancel gravity.  In simulation the physics engine
handles this, so gravity compensation can be disabled.
\end{reasonbox}


% ══════════════════════════════════════════════════════════════════════════════
\section{Reference Generation and Interpolation}
\label{sec:refgen}

The controller receives targets asynchronously at a lower frequency and must
output smooth references at the control rate.

% --------------------------------------------------------------------------
\subsection{Position Mode --- SLERP and Linear Interpolation}
\label{sec:pos_interp}

When $t_{\text{rem}} > 0$ (remaining time to target):

\paragraph{Translation} (per-cycle linear interpolation):
\begin{equation}
  p_{\text{ref}}(k\!+\!1)
  = p_{\text{ref}}(k)
    + \frac{p_{\text{target}} - p_{\text{ref}}(k)}{f_c \cdot t_{\text{rem}}}
  \label{eq:lin_interp}
\end{equation}
where $f_c$ is the control frequency (Hz).

\paragraph{Rotation} (SLERP --- Spherical Linear Interpolation):
\begin{equation}
  q_{\text{ref}}(k\!+\!1)
  = \SLERP\!\!\left(q_{\text{ref}}(k),\; q_{\text{target}},\;
                     \frac{1}{f_c \cdot t_{\text{rem}}}\right)
  \label{eq:slerp}
\end{equation}

SLERP interpolates along the great arc on the unit quaternion sphere:
\[
  \SLERP(q_0,\,q_1,\,t)
  = \frac{\sin\bigl((1-t)\,\Omega\bigr)}{\sin\Omega}\,q_0
    + \frac{\sin(t\,\Omega)}{\sin\Omega}\,q_1,
  \qquad
  \Omega = \arccos(q_0 \cdot q_1)
\]
ensuring constant angular velocity.

When $t_{\text{rem}} \le 0$ (the common case): the reference \emph{snaps} to
the target and the reference velocity is set to zero.

% --------------------------------------------------------------------------
\subsection{Velocity Mode --- Lie-Group Integration}
\label{sec:vel_interp}

In velocity mode, the target provides a twist $\bm{\xi}_d$ in the TCP frame.
The reference pose is integrated on $\SE(3)$:

\begin{keqbox}[Body-Frame Velocity Integration]
\begin{equation}
  T_{\text{ref}}(k\!+\!1)
  = T_{\text{ref}}(k) \cdot \Exp(\bm{\xi}_d\,\Delta t)
  \label{eq:vel_integ}
\end{equation}
\end{keqbox}

where $\Delta t = 1/f_c$.

\begin{reasonbox}[Why right-multiply?]
The twist $\bm{\xi}_d$ is in the TCP (body) frame, so we right-multiply.
This means ``move forward along my current heading'' regardless of the global
orientation --- exactly the right semantics for teleop or tool-frame velocity
commands.
\end{reasonbox}

After integration, the pose is clamped to Cartesian limits (\S\ref{sec:clamping}).

For joint velocity mode:
\[
  q_{\text{ref}}(k\!+\!1) = q_{\text{ref}}(k) + \dot{q}_d\,\Delta t
\]


% ══════════════════════════════════════════════════════════════════════════════
\section{Cartesian and Joint Limit Clamping}
\label{sec:clamping}

% --------------------------------------------------------------------------
\subsection{Translation Clamping}

Position targets are clamped element-wise to a bounding box:
\[
  p_i = \clamp(p_i,\; p_{\min,i},\; p_{\max,i}),
  \qquad i \in \{x,\,y,\,z\}
\]

% --------------------------------------------------------------------------
\subsection{Rotation Clamping via Quaternion Log-Map}
\label{sec:rot_clamp}

Rotation limits are expressed in the tangent space relative to a reference
quaternion $q_{\text{ref}}$:

\begin{enumerate}[nosep]
  \item Relative quaternion:
    $q_{\text{rel}} = q_{\text{target}} \cdot q_{\text{ref}}^{-1}$
  \item Log-map:
    $\bm{\phi} = \Log(q_{\text{rel}}) \in \R^3$
  \item Clamp:
    $\phi_i = \clamp(\phi_i,\;\phi_{\min,i},\;\phi_{\max,i})$
  \item Reconstruct:
    $q_{\text{clamped}} = \Exp(\bm{\phi}_{\text{clamped}}) \cdot q_{\text{ref}}$
\end{enumerate}

\begin{reasonbox}[Why log-map instead of Euler angles?]
The log-map provides a singularity-free, minimal parameterisation of the
relative rotation.  Euler angles suffer from gimbal lock and ambiguous
representations near $\pm 90^\circ$ pitch.
\end{reasonbox}

% --------------------------------------------------------------------------
\subsection{Velocity Scaling}

Velocity targets are \emph{scaled} (not clamped) to preserve direction:
\[
  \dot{p}_{\text{scaled}} = s\,\dot{p},
  \qquad
  s = \min_i\!\left(\frac{v_{\max,i}}{|\dot{p}_i|}\right) \le 1
\]

A single scaling factor across all translational axes preserves direction.
For rotational velocity:
\[
  \|\bm{\omega}\| > \omega_{\max}
  \;\;\Longrightarrow\;\;
  \bm{\omega}_{\text{scaled}} = \omega_{\max} \cdot \frac{\bm{\omega}}{\|\bm{\omega}\|}
\]

% --------------------------------------------------------------------------
\subsection{Bisquare Soft-Margin Deceleration}
\label{sec:bisquare}

When a soft margin $m > 0$ is configured, velocity is smoothly scaled to zero
as position approaches a limit.  The normalised penetration is:
\[
  d = \clamp\!\left(\frac{\text{distance from margin boundary}}{m},\;0,\;1\right)
\]

The scale factor uses the \textbf{bisquare} (biweight) kernel:

\begin{keqbox}[Bisquare Deceleration]
\begin{equation}
  s(d) = (1 - d^2)^2
  \label{eq:bisquare}
\end{equation}
\end{keqbox}

Properties:
\begin{itemize}[nosep]
  \item $s(0) = 1$: full speed at the margin boundary
  \item $s(1) = 0$: zero speed at the hard limit
  \item $s'(0) = 0$, $s'(1) = 0$: smooth onset/offset --- no velocity discontinuities
\end{itemize}

\begin{center}
\begin{tikzpicture}[scale=3.5, >=Stealth, font=\small]
  \draw[->] (-0.05,0) -- (1.15,0) node[right]{$d$};
  \draw[->] (0,-0.05) -- (0,1.15) node[above]{$s(d)$};
  \draw[domain=0:1, samples=100, thick, blue!70!black]
    plot (\x, {(1-\x*\x)*(1-\x*\x)});
  \foreach \x/\lab in {0/0, 0.5/0.5, 1/1}{
    \draw (\x,0.02) -- (\x,-0.02) node[below]{\lab};
  }
  \foreach \y/\lab in {0.5/0.5, 1/1}{
    \draw (0.02,\y) -- (-0.02,\y) node[left]{\lab};
  }
  \node[blue!70!black, above right] at (0.35, 0.6) {$(1-d^2)^2$};
\end{tikzpicture}
\end{center}


% ══════════════════════════════════════════════════════════════════════════════
\section{Impedance Parameter Smoothing}
\label{sec:smoothing}

% --------------------------------------------------------------------------
\subsection{Exponential Smoothing of Stiffness and Damping}

Stiffness and damping are not applied instantaneously.  An exponential moving
average smooths transitions:

\begin{keqbox}[Exponential Smoothing]
\begin{equation}
  S_{n+1} = (1 - c)\,S_n + c\,S_{\text{target}}
  \label{eq:ema}
\end{equation}
\end{keqbox}

where $c \in [0,1]$ is the smoothing constant (separate values for stiffness
and damping).

\begin{reasonbox}
Instantaneous stiffness changes produce torque discontinuities that can excite
structural resonances or jerk the arm.  Smoothing acts as a first-order
low-pass filter with effective time constant $\tau \approx \Delta t / c$.
\end{reasonbox}

\begin{examplebox}
With $c = 0.05$ at 1\,kHz, the time constant is
$\tau = 1\,\text{ms}/0.05 = 20$\,ms.  A step in stiffness from 100 to
500\,N/m reaches 95\% in $\approx 60$\,ms ($\approx 3\tau$).
\end{examplebox}

% --------------------------------------------------------------------------
\subsection{Feedforward Wrench Interpolation (Rate-Limited)}

The feedforward wrench is ramped toward its target using a rate limiter:
\begin{equation}
  F_{\text{ff},i}(k\!+\!1)
  = F_{\text{ff},i}(k)
    + \clamp\!\bigl(
        F_{\text{target},i} - F_{\text{ff},i}(k),\;
        -\dot{F}_{\max,i}\,\Delta t,\;
         \dot{F}_{\max,i}\,\Delta t
      \bigr)
  \label{eq:rate_limit}
\end{equation}

The target wrench is also clamped to configurable min/max bounds before ramping.

% --------------------------------------------------------------------------
\subsection{Force Feedback with Tare Offset}
\label{sec:force_fb}

A feedback loop mixes the feedforward wrench with the sensed wrench:
\begin{equation}
  F_{\text{total}}
  = F_{\text{ff}}
    + G_{\text{fb}} \odot \bigl(F_{\text{ff}} - F_{\text{sensed}}^{\text{tared}}\bigr)
  \label{eq:force_fb}
\end{equation}

where $G_{\text{fb}} \in \R^6$ are per-axis feedback gains and
$F_{\text{sensed}}^{\text{tared}} = F_{\text{sensed}} - F_{\text{tare}}$.

The tare offset is captured by a service call and stored in the base frame.
Each cycle it is transformed to the TCP frame:
\[
  F_{\text{tare}}^{\text{tip}} = R_{\text{tcp}}^{-1}\,F_{\text{tare}}^{\text{base}}
\]

The total wrench (TCP frame) is then transformed to base frame for the
impedance law:
\[
  F_{\text{ff}}^{\text{base}} = R_{\text{tcp}}\,F_{\text{total}}^{\text{tcp}}
\]
(Applied separately to force and torque 3-vectors.)


% ══════════════════════════════════════════════════════════════════════════════
\section{Tracking Error Watchdog}
\label{sec:watchdog}

A safety mechanism detects when the controller cannot make progress:

\begin{enumerate}[nosep]
  \item Each cycle, compute
    $\Delta e = |e_x(k) - e_x(k\!-\!1)|$.
  \item If $\|e_x\|_\infty >$ minimum error threshold \textbf{and}
    $\Delta e$ is below minimum change threshold for longer than a configurable
    timeout $\Longrightarrow$ reset target to current pose.
\end{enumerate}

This prevents indefinite build-up of integral error or feedforward wrench when
blocked by obstacles, joint limits, or singularities.


% ══════════════════════════════════════════════════════════════════════════════
\section{Complete Algorithm Flow}
\label{sec:flow}

\begin{center}
\begin{tikzpicture}[
  node distance=8mm and 5mm,
  block/.style={draw, rounded corners=3pt, fill=boxblue,
                minimum width=5.8cm, minimum height=8mm,
                font=\small\sffamily, align=center},
  arrow/.style={->, thick, >=Stealth},
  font=\small\sffamily
]

\node[block] (read) {1. Read Hardware\\
  {\scriptsize $q,\,\dot{q}$ from joints;\; FK $\to$ $T,\,\xi$;\; FT sensor}};

\node[block, below=of read] (ingest) {2. Ingest Commands\\
  {\scriptsize MotionUpdate / JointMotionUpdate}};

\node[block, below=of ingest] (ref) {3. Generate Reference\\
  {\scriptsize Clamp to limits;\; SLERP / linear interp / velocity integration}};

\node[block, below=of ref] (smooth) {4. Smooth Impedance Parameters\\
  {\scriptsize EMA on $K,D$;\; rate-limit $F_{\text{ff}}$;\; force feedback + tare}};

\node[block, below=of smooth] (compute) {5. Compute Control Torques\\
  {\scriptsize Cartesian: $\tau\!=\!J^\tr F + \tau_{\text{ns}} + \tau_{\text{avoid}}$
  \;\;or\;\; Joint: $\tau\!=\!K \odot e_q + D \odot e_{\dot{q}} + \tau_{\text{ff}}$
  \;\;$+\;\tau_g$}};

\node[block, below=of compute] (write) {6. Write Hardware\\
  {\scriptsize Set effort command interfaces}};

\node[block, below=of write] (pub) {7. Publish State\\
  {\scriptsize TCP pose/vel, error, joint ref, tare offset}};

\draw[arrow] (read) -- (ingest);
\draw[arrow] (ingest) -- (ref);
\draw[arrow] (ref) -- (smooth);
\draw[arrow] (smooth) -- (compute);
\draw[arrow] (compute) -- (write);
\draw[arrow] (write) -- (pub);

\end{tikzpicture}
\end{center}

\subsection{Final Torque Equations}

\paragraph{Cartesian impedance mode:}

\begin{keqbox}[Complete Cartesian Impedance Torque]
\begin{equation}
\begin{aligned}
  \tau \;=\;&
    \underbrace{
      J^\tr\!\bigl(
        K\,e_x + D\,e_v + F_{\text{ff}}
        + K_I \odot e_I + F_{\text{offset}}
      \bigr)
    }_{\mathclap{\text{Cartesian impedance}}}
  \\\[6pt]
  &+\; \underbrace{
    (I - JJ^\dagger)\,
    \bigl(K_{\text{ns}} \odot (q^*_{\text{ns}} \!-\! q)
          - D_{\text{ns}} \odot \dot{q}\bigr)
  }_{\mathclap{\text{Nullspace}}}
  \\\[6pt]
  &+\; \underbrace{\tau_{\text{avoid}}}_{\mathclap{\text{Joint limit}}}
  \;+\; \underbrace{\tau_g(q)}_{\mathclap{\text{Gravity}}}
\end{aligned}
\label{eq:final_cart}
\end{equation}
\end{keqbox}

\paragraph{Joint impedance mode:}

\begin{keqbox}[Complete Joint Impedance Torque]
\begin{equation}
\begin{aligned}
  \tau \;=\;&
    K \odot (q^* \!-\! q)
    + D \odot (\dot{q}^* \!-\! \dot{q})
  \\\[4pt]
  &+\; \tau_{\text{ff}} + \tau_g(q)
\end{aligned}
\label{eq:final_joint}
\end{equation}
\end{keqbox}


% ══════════════════════════════════════════════════════════════════════════════
\section{File-to-Concept Map}

\begin{center}
\small
\begin{tabular}{@{}p{6.7cm}p{7.7cm}@{}}
\toprule
\textbf{File} & \textbf{Sections} \\
\midrule
\path{src/actions/cartesian_impedance_action.cpp}
  & \S\ref{sec:cartesian_imp} (Wrench, $J^\tr$),
    \S\ref{sec:nullspace} (Nullspace),
    \S\ref{sec:jlimit} (Joint limits) \\[4pt]
\path{include/.../cartesian_impedance_action.hpp}
  & \S\ref{sec:svd} (Chiaverini pseudo-inverse docs) \\[4pt]
\path{src/actions/joint_impedance_action.cpp}
  & \S\ref{sec:joint_imp} (Joint impedance) \\[4pt]
\path{src/actions/gravity_compensation_action.cpp}
  & \S\ref{sec:gravity} (KDL gravity) \\[4pt]
\path{src/aic_controller.cpp}
  & \S\ref{sec:refgen} (Reference gen),
    \S\ref{sec:clamping} (Clamping),
    \S\ref{sec:smoothing} (Smoothing),
    \S\ref{sec:watchdog} (Watchdog),
    \S\ref{sec:flow} (Full flow) \\[4pt]
\path{src/utils.cpp}
  & \S\ref{sec:explog} (Exp/Log),
    \S\ref{sec:vel_interp} (SE3 integration) \\[4pt]
\path{include/.../cartesian_limits.hpp}
  & \S\ref{sec:rot_clamp} (Rotation limits) \\[4pt]
\path{include/.../cartesian_state.hpp}
  & \S\ref{sec:SE3} (Pose representation) \\
\bottomrule
\end{tabular}
\normalsize
\end{center}


% ══════════════════════════════════════════════════════════════════════════════
\section*{References}

\begin{enumerate}[label={[\arabic*]}, nosep]
  \item Chiaverini, S. (1997).
    \emph{Singularity-robust task-priority redundancy resolution for real-time
    kinematic control of robot manipulators.}
    IEEE Transactions on Robotics and Automation, 13(3), 398--410.

  \item Siciliano, B., Sciavicco, L., Villani, L., \& Oriolo, G. (2009).
    \emph{Robotics: Modelling, Planning and Control.}
    Springer.

  \item Lynch, K.\,M. \& Park, F.\,C. (2017).
    \emph{Modern Robotics: Mechanics, Planning, and Control.}
    Cambridge University Press.
\end{enumerate}

\end{document}
